\section{One-loop normalization and trace reflection}\label{sec:loop-normal}

\subsection{Certified iteration}

A loop node is specified by:
\begin{itemize}
  \item a packed state code \(s\);
  \item an initializer \(s_0(\vec x)\);
  \item a transition \(T(i,s,\vec x)\);
  \item a fuel term \(F(\vec x)\);
  \item a quantifier-free invariant \(I(i,s,\vec x)\);
  \item an output projection \(O(s)\).
\end{itemize}
All numerical maps are primitive recursive.  The state trace is the primitive recursion
\[
 S(0,\vec x)=s_0(\vec x),
 \qquad
 S(i+1,\vec x)=T(i,S(i,\vec x),\vec x).
\]

\begin{theorem}[Object trace reflection]\label{thm:trace-reflection}
Assume object-level derivations of
\begin{align*}
 &I(0,s_0(\vec x),\vec x),\\
 &I(i,s,\vec x)\longrightarrow I(i+1,T(i,s,\vec x),\vec x).
\end{align*}
Then \iPRA\ derives \(I(i,S(i,\vec x),\vec x)\) for all open \(i\), and in particular at the fuel term \(F(\vec x)\).  If a terminal derivation proves
\[
 I(F,S(F),\vec x)\longrightarrow Q(O(S(F)),\vec x),
\]
then \(Q\) is derivable.
\end{theorem}

\begin{proof}
Substitute the trace term into the invariant to obtain the QF formula
\[
 J(i,\vec x)\equiv I(i,S(i,\vec x),\vec x).
\]
The base derivation is the supplied base proof followed by the zero computation equation for \(S\).  For the step, assume \(J(i,\vec x)\), use the supplied transition-preservation proof, and rewrite \(S(i+1)\) by the successor computation equation.  Apply quantifier-free induction to \(J\).  Substitution instantiates the resulting open theorem at \(F(\vec x)\), and modus ponens with the terminal derivation yields \(Q\).
\end{proof}

This theorem is already very close to the current \texttt{LeanProofs.IPRA.ind} API.  The accompanying \texttt{KernelCore.lean} packages exactly the base, step, and QF fields needed to invoke it.

\subsection{Why all precision must be fixed before induction}

A common invalid pattern is to put into the invariant the assertion
\[
  \forall q\; \text{``the state is correct to precision }q\text{''}.
\]
This is not quantifier-free.  Even when every individual precision computation is finite, the displayed invariant is stronger than the induction rule permits.

The correct transformation is:
\begin{enumerate}[label=(\arabic*)]
  \item fix the terminal precision \(q_*\) as an open parameter;
  \item run the reverse demand pass to calculate every intermediate request from \(q_*\);
  \item allocate in the state only the finite observations required at those requests;
  \item prove one QF invariant about this finite table;
  \item after obtaining the open theorem for arbitrary \(q_*\), universally quantify over \(q_*\) if the final theorem needs all precisions.
\end{enumerate}
The universal quantifier is introduced after induction, not inside its matrix.

\subsection{Collapsing nested bounded constructions}

A fixed theorem may contain many loops, searches, and finite table passes.  They can be combined into one state recursion.

\begin{theorem}[One-loop normal form]\label{thm:one-loop}
Every fixed program generated from primitive-recursive assignments, finite conditionals, bounded loops, bounded searches, finite-table folds, and acyclic calls to previously generated routines can be transformed into:
\begin{enumerate}[label=(\roman*)]
  \item one primitive-recursive fuel term \(F(\vec x)\);
  \item one packed state \(s\) containing a program counter, local registers, and loop counters;
  \item one primitive-recursive transition \(T\);
  \item primitive-recursive projections recovering the original outputs.
\end{enumerate}
If every source loop has a QF invariant and every source command has a local preservation derivation, the combined machine has a QF invariant and an object-level base and step derivation.
\end{theorem}

\begin{proof}
Because the source program is fixed, its control-flow graph has finitely many labels.  Allocate a numerical program counter \(pc<L\).  Allocate a fixed tuple of registers large enough to hold every live local variable; registers not live at the current label are ignored.  Replace each atomic command by one branch of a finite selector on \(pc\).  A loop header tests its counter and jumps either to the body or the successor label.  A bounded search is its explicit \((found,witness,index)\) loop.  A finite table fold is an index/accumulator loop.

Define a structural cost function in the metatheory.  An atomic command costs one step; sequential costs add; a loop with bound \(B\) costs a primitive-recursive expression obtained from \(B\) and a bound on the body cost; fixed nesting preserves primitive recursiveness.  Add a small constant for control transfers.  This gives total fuel \(F\).  Alternatively, include a halting flag and allow idle transitions after the source program halts; then any certified upper bound on the runtime is sufficient.

The combined invariant is a finite conjunction indexed by control labels:
\[
  \bigwedge_{\ell<L}\bigl(pc=\ell\to I_\ell(\text{live registers})\bigr),
\]
together with range equations for the program counter, flags, and bounded indices.  Since \(L\) is fixed and every \(I_\ell\) is QF, the conjunction is QF.  The base proof establishes the invariant at the entry label.  The step proof performs a finite case split on the reflected value of \(pc\), invokes the local preservation derivation for that command, and establishes the invariant at the successor label.  Apply \cref{thm:trace-reflection}.  Output projection equations show that the packed final state agrees with the source program.
\end{proof}
\subsection{Clocked traces as an alternative}
Sometimes it is easier to let an external producer supply a trace and have \iPRA\ verify it.  A trace certificate contains \(s_0,\ldots,s_F\).  A primitive-recursive checker verifies the initial state, every adjacent transition, and the terminal predicate.  Boolean fold reflection turns checker acceptance into the conjunction of local facts.
This route is useful when the producer is complex but the transition is simple.  However, the checker still needs an object-level reflection theorem, and the trace length must have an explicit bound.  ``Python produced a valid trace'' is not itself the desired derivation; the checked acceptance equation and its reflection proof are.
\subsection{Lexicographic descent and fuel}
Algorithms are often justified by a decreasing tuple rather than an obvious loop bound.  Let
\[
 \rho(s)=(r_0(s),\ldots,r_{k-1}(s))
\]
take values in a bounded box \(0\le r_j<M_j\), and assume every nonterminal transition decreases \(\rho\) lexicographically.  Encode the tuple by the mixed-radix rank
\[
 R(\rho)=r_0\prod_{j>0}M_j+r_1\prod_{j>1}M_j+\cdots+r_{k-1}.
\]
Lexicographic decrease implies strict numerical decrease of \(R\).  Therefore the initial rank plus one is a primitive-recursive fuel bound.  The proof uses only finite arithmetic and is suitable for the algebraic backend.  This resolves many ``well-founded recursion'' uses in the packets without introducing transfinite induction.
